%HOMEWORK template for M 325K
\documentclass{article}
\headheight=7pt         \topmargin=14pt
\textheight=574pt       \textwidth=445pt
\oddsidemargin=18pt     \evensidemargin=18pt

%Begin package import

\usepackage{amsmath, amssymb, amsthm}
\usepackage{mathtools}
\usepackage{pdfpages}
\usepackage{tikz-cd}

%End package import
%Begin custom commands

\newcommand{\C}{\mathbb{C}}
\newcommand{\R}{\mathbb{R}}
\newcommand{\Q}{\mathbb{Q}}
\newcommand{\Z}{\mathbb{Z}}
\newcommand{\N}{\mathbb{N}}
\newcommand{\F}{\mathbb{F}}
\newcommand{\abs}[1]{\left\lvert #1\right\rvert}
\newcommand{\RM}[1]{\mathrm{#1}}
\newcommand{\BF}[1]{\textbf{#1}}
\newcommand{\e}{\epsilon}
\newcommand{\ceil}[1]{\lceil{#1}\rceil}
\newcommand{\floor}[1]{\lfloor{#1}\rfloor}
\newcommand{\rarr}[0]{\rightarrow}
\newcommand{\IM}[1]{\text{Im}(#1)}
\newcommand{\Hom}[0]{\text{Hom}}
\newcommand{\Pow}[1]{\text{Pow}(#1)}
\newcommand{\angles}[1]{\langle #1 \rangle}

%End custom commands
%Begin custom theorems

\newtheorem{innercustomgeneric}{\customgenericname}
\providecommand{\customgenericname}{}
\newcommand{\newcustomtheorem}[2]{%
\newenvironment{#1}[1]
{%
\renewcommand\customgenericname{#2}%
\renewcommand\theinnercustomgeneric{##1}%
\innercustomgeneric%
}
{\endinnercustomgeneric}
}

\newcustomtheorem{THRM}{Theorem}
\newcustomtheorem{LEM}{Lemma}
\newcustomtheorem{EX}{Exercise}
\newcustomtheorem{COR}{Corollary}
\newcustomtheorem{PROB}{Problem}
\newcustomtheorem{claim}{Claim}

\newenvironment{solution}
{\renewcommand\qedsymbol{$\blacksquare$}\begin{proof}[Solution]}
{\end{proof}}

\newenvironment{definition}
{\renewcommand\qedsymbol{$\blacksquare$}\begin{proof}[Definition]}
{\end{proof}}

%End custom theorems
\begin{document}

\title{Homework 9}
\author{Arham Lodha}%put your name here
\date{March 26, 2024}%enter the due date here
\maketitle

\begin{PROB}{1}\label{Problem 1.}
\end{PROB}
\begin{proof}
    Let $(x_n)$ be a sequence of positive terms which converge to $x > 0$. Convergent sequences are bounded, so $\exists M \in \R \ni x_n \leq M$. $0 < x_n < M \implies x_n = \abs{x_n} \leq M$. $\forall \epsilon > 0$. By the definition of limit, $\exists K \in \N \ni \forall n \geq K, \abs{x_n - x} < \frac{\epsilon}{\abs{x} + \abs{M}}$.     
    
    \begin{align*}
        \abs{x_n^2 - x^2} &= \abs{x_n^2 - x x_n + x x_n - x^2} \\
        &= \abs{(x_n - x)(x + x_n)} \\
        &= \abs{x_n - x}\abs{x + x_n} \\
        &\leq \abs{x_n - x}(\abs{x} + \abs{x_n}) \text{ (triangle inequality)} \\
        &\leq \abs{x_n - x}(\abs{x} + \abs{M}) \text{ (since $\abs{x_n} \leq M$)} \\
        &< \frac{\epsilon}{\abs{x} + \abs{M}}(\abs{x} + \abs{M}) \\
        &= \epsilon.
    \end{align*}
    
    Thus $\lim x_n^2 = x^2$ 
\end{proof}

\bigskip

\begin{PROB}{2}\label{Problem 2.}
\end{PROB}
\begin{proof}
    $\forall n \in \N$, $x_n = \frac{1}{\frac{1}{n} + 4} = \frac{n}{1 + 4n}$.
    
    $\forall \epsilon > 0$. By Archemedian Property, $\exists K \in \N \ni K > \frac{1}{16\epsilon} \implies \epsilon > \frac{1}{16K}$. $\forall n \geq K$   
    
    
    \begin{align*}
        \abs{\frac{n}{4n + 1} - \frac{1}{4}} &= \abs{\frac{4n - (4n + 1)}{16n + 4}} \\
        &= \abs{\frac{-1}{16n + 4}} \\
        &= \frac{1}{16n + 4} \\
        &< \frac{1}{16n} \text{ (since $16n + 4 > 16n$ )}\\
        &\leq \frac{1}{16K} \text{ (since $n \geq K \implies 16n \geq K$ )}\\
        &< \epsilon
    \end{align*}

    Thus, $\lim x_n = \frac{1}{4}$. 
\end{proof}

\bigskip

\begin{PROB}{3a}\label{Problem 3a.}
\end{PROB}
\begin{proof}
    Suppose $x_n = \frac{(-1)^n}{n^3}$. 
    
    $\forall \epsilon > 0$. By Archemedian property, $\exists K \in \N \ni K \geq \frac{1}{\epsilon^{\frac{1}{3}}} \implies \epsilon > \frac{1}{K^3}$. $\forall n \geq K$. 
    
    \begin{align*}
        \abs{\frac{(-1)^n}{n^3}} &= \frac{1}{n^3} \\ 
        &\leq \frac{1}{K^3} \text{ (since $n \geq K \implies n^3 \geq K^3$)} \\
        &< \epsilon \text{ (since $\epsilon > \frac{1}{K^3}$)}
    \end{align*}

    Thus, $\lim x_n = 0$ 
\end{proof}

\bigskip

\begin{PROB}{3b}\label{Problem 3b.}
\end{PROB}
\begin{proof}
    Suppose $x_n = (-1)^n(1 + \frac{1}{n})$. 

    Take the subsequences, $y_n = x_{2n} = (-1)^{2n}(1 + \frac{1}{2n}) = 1 + \frac{1}{2n}$ and $z_n = x_{2n - 1} = (-1)^{2n - 1}(1 + \frac{1}{2n - 1}) = -1 -\frac{1}{2n - 1}$.

    Limit of constants are equal to the constant, thus $\lim_{n \to \infty} 1 = 1$ and $\lim_{n \to \infty} (-1) = -1$.
    
    
    $\forall \epsilon > 0$. By Archemedian property, $\exists K \in \N \ni K \geq \frac{1}{2\epsilon} + 0.5 \implies \epsilon > \frac{1}{2K - 1}$. $\forall n \geq K$.

    \begin{align*}
        \abs{\frac{1}{2n - 1}} &= \frac{1}{2n - 1} \\
        &\leq \frac{1}{2K - 1} \text{ ( since $n \geq K \implies 2n - 1 \geq 2K - 1$)} \\
        &< \epsilon 
    \end{align*}

    Thus $\lim \frac{1}{2n - 1} = 0$. 
    
    By operation of limits theorem, $\lim_{n \to \infty} y_n = \lim_{n \to \infty}(1 + \frac{1}{2n}) = \lim_{n \to \infty}(1) + \frac{1}{2}\lim_{n \to \infty}\frac{1}{n} = 1$ and $\lim_{n \to \infty} z_n = \lim_{n \to \infty}(-1 - \frac{1}{2n - 1}) = \lim_{n \to \infty}(-1)- \lim_{n \to \infty}(\frac{1}{2n - 1}) = -1 - 0 = -1$.
    
    Thus $\lim y_n = 1$ and $\lim z_n = -1$. Thus subsequence of $x_n$ have different limits, thus because of the divergence criteria given by subsequences, $x_n$ is divergent.      
\end{proof}

\begin{PROB}{3c}\label{Problem 3c.}
\end{PROB}
\begin{proof}
    Suppose $x_n = (-1)^n n^3$. 

    $\forall r \in \R$. 
    
    \begin{enumerate}
        \item If $r \leq 0$. $x_1 = 1 > r \implies r$ is not a upper bound for $(x_n)$.
        \item If $r > 0$, 
        \begin{enumerate}
            \item If $\ceil{r}$ is even, let $K := \ceil{r} + 2$, thus $K$ even and $K > \ceil{r} \geq r$. Note $K > 1$. Thus $x_K = (-1)^K K^3 = K^3 > K > r$.
            \item If $\ceil{r}$ is odd, let $K = \ceil{r} + 1$, thus $K$ even and $K > \ceil{r} \geq r$. Note $K > 1$. Thus $x_n = (-1)^K K^3 = K^3 > K>r$.      
        \end{enumerate}

        Thus $r$ is not a upper bound for $(x_n)$.   
    \end{enumerate}

    Thus $(x_n)$ is unbounded. By contrapositive of the convergence $\implies$ bounded theorem, $(x_n)$ is divergent.
\end{proof}

\begin{PROB}{3d}\label{Problem 3d.}
\end{PROB}
\begin{proof}
    Suppose $x_n = \frac{3n^3 - n + 4}{6n^3 + n + 5}$. $\forall \epsilon > 0$, By the Archemedian property, $\exists K \in \N \ni K > \frac{\sqrt{3}}{2\sqrt{\epsilon}} \implies \epsilon > \frac{3}{14K^2}$     

    
    \begin{align*}
        \abs{\frac{3n^3 - n + 4}{6n^3 + n + 5} - \frac{1}{2}} &= \abs{\frac{6n^3 - 2n + 8 - (6n^3 + n + 5)}{12n^3 + 2n + 10}} \\
        &= \abs{\frac{-3n + 3}{2(6n^3 + n + 5)}} \\
        &= \abs{\frac{-3}{2} * \frac{n - 1}{6n^3 + n + 5}}\\
        &= \frac{3}{2}\left(\frac{n - 1}{6n^3 + n + 5}\right) \\
        &< \frac{3}{2}\left(\frac{n}{6n^3 + n + 5}\right) \\
        &< \frac{3}{2}\left(\frac{n}{6n^3 + n}\right) \\
        &< \frac{3}{2}\left(\frac{1}{6n^2 + n^2}\right) \\
        &< \frac{3}{2}\left(\frac{1}{7n^2}\right) \\
        &\leq \frac{3}{14K^2} \\
        &< \epsilon
    \end{align*}

    Thus $\lim x_n = \frac{1}{2}$. 
\end{proof}

\bigskip

\begin{PROB}{4a}\label{Problem 4a.}
\end{PROB}
\begin{proof}
    (Base Case): For $n = 0$, $x_0 = 5 > -3$

    (Inductive Hypothesis): Suppose the statement holds for $n $ ie suppose $x_{n} > -3$

    (Inductive step): $x_{n} > -3 \implies x_{n} / 3 > -1 \implies x_{n}/3 - 2 = x_{n + 1} > -1 - 2 = -3$. 
    
    Thus by induction, $x_n > -3$. 
\end{proof}
\begin{PROB}{4b}\label{Problem 4b.}
\end{PROB}
\begin{proof}
    (Base Case): For $n = 0$, $x_0 = 5$ and $x_1 = \frac{5}{3} - 2 = -\frac{1}{3} < x_0$. Thus $x_0 < x_1$
    
    (Inductive Hypothesis): Suppose $x_{n} < x_{n + 1}$
    
    (Inductive Step): $x_{n} < x_{n + 1} \implies x_{n} / 3 < x_{n + 1} / 3 \implies x_{n} / 3 - 2 = x_{n + 1} < x_{n + 1} /3 - 2 = x_{n + 2}$.
    
    Thus by induction, $(x_n)$ is monotonically decreasing.  
\end{proof}
\begin{PROB}{4c}\label{Problem 4c.}
     
\end{PROB}
\begin{proof}
    By monotonic convergence theorem, since $(x_n)$ has a lower bound and is decreasing, $(x_n)$ is convergent.  
\end{proof}
\begin{PROB}{4d}\label{Problem 4d.}
\end{PROB}
\begin{proof}
    Let $L := \lim x_n$. 
    
    $\lim x_{n + 1} = L$ since limit of a subsequence is equal to the limit of the sequence. 
    
    By the operation of limits, $\lim x_{n + 1} = \frac{1}{3}\lim x_n + \lim (-2)$. Limit of a constant term, is the constant, thus $\lim x_{n + 1}= \frac{1}{3}L - 2$. 
    
    By Uniqueness of Limits, $\frac{1}{3}L - 2 = L \implies -2 = \frac{2}{3}L \implies L = -3$.  
\end{proof}

\bigskip
\begin{PROB}{5a}\label{Problem 5a.}
\end{PROB}
\begin{proof}
    $\forall \alpha \in \R$. By the definition of proper divergence, $\exists K \in \N \ni \forall n \geq K, x_n > \alpha \implies -x_n < -\alpha$, since multiplication by a negative real number flips the ordering sign. Thus $-x_n < -\alpha$. 
    
    By the definition of proper divergence, $(-x_n)$ diverges to $- \infty$.  
\end{proof}
\begin{PROB}{5b}\label{Problem 5b.}  
\end{PROB}
\begin{proof}
    Let $\alpha = 1$. For this $\alpha$, $\exists K \ni \forall n \geq K, x_n > \alpha$. $\forall m, n \in \N \ni m > n \geq K$. 
    
    \begin{align*}
        \abs{\sum_{k = 1}^{m}x_k - \sum_{k = 1}^{n}x_k} &= \abs{\sum_{k = n + 1}^{m}x_k} \\
        &= \sum_{k = n + 1}^{m}x_k \text{ (since $x_k > \alpha > 0$ $\forall k \in \N \ni n < k \leq m$)} \\
        &> (m - n) \text{(since $x_k > 1$)} \\
        &> 1
    \end{align*}

    Thus by definition of cauchy, $(\sum x_n)$ is not cauchy. Since convergence $\iff$ cauchy, $(\sum x_n)$ is not diverges.  
\end{proof}

\end{document}